% GENERATED FILE. Edit canonical lessons, then run npm run generate:books.
% canonical-source-hash: fnv1a64:2cf4776e27657fe2
% canonical-lessons: ES-C442-administracion, ES-C442-envio, ES-C442-concertar, ES-C442-reunion, ES-C442-pantalla, ES-R442-repaso-gestion, ES-C442-sintesis-gestion

\chapter{La Gestion}
\label{ch:es-gestion}
\hlchaptermodality{442}

\begin{chapteropening}
\textbf{By the end of this chapter:} \emph{I can deal with an institution rather than a person: book the appointment, read what the office wants, track what it sent me, and do it on a screen.}

Read the message or notice this chapter's words exist for, then say the same thing back.
\end{chapteropening}

\section[la administración]{la administración — the lesser one, serving}

\label{lesson:ES-C442-administracion}



\begin{quote}
Say \textbf{la oficina}, \textbf{el documento} and \textbf{servir}. Put the three together and you have named the thing that asks you for papers.
\end{quote}

\subsection*{You'll want to know: la administración}
\textbf{la administración} — the administration, and in Spain usually \textbf{the state}: \emph{la Administración} with a capital letter is officialdom as a whole.

\emph{Trabaja en la administración} — he works in the public sector. \emph{Trámites con la administración} — dealings with officialdom. The verb is \textbf{administrar}, to manage or to administer.

The ending is \textbf{-ción}, which you have had since the chapter on paperwork, so the word is \emph{the act of administering} made into a thing.

\begin{cousinweb}
\textbf{ad-} towards, plus \textbf{ministr-} from Latin \textbf{minister}, a servant or attendant, plus \textbf{-ción}.

And \emph{minister} comes from \textbf{minus}, less. A \emph{minister} was \textbf{the lesser one}, the person who serves — which is why the word for a government's most senior officials is built on the word for \emph{less}.

The pair is the point:

\noindent
\begin{tabularx}{\linewidth}{@{}>{\raggedright\arraybackslash}X>{\raggedright\arraybackslash}X>{\raggedright\arraybackslash}X@{}}
\toprule
\textbf{Latin word} & \textbf{built on} & \textbf{what it named} \\
\midrule
\textbf{minister} & \emph{minus}, less & the one who serves \\
\textbf{magister} & \emph{magis}, more & the one who directs \\
\bottomrule
\end{tabularx}

\emph{Magister} is where \textbf{el maestro} comes from. So \emph{maestro} and \emph{ministro} are the same shape twice over, once on \emph{more} and once on \emph{less}, and Spanish kept both.

You already own the smaller half of that pair outright: \textbf{menos} is \emph{minus}, unchanged in meaning and barely changed in sound. The administration is named with the same word you use to say \emph{less}.
\end{cousinweb}

\begin{culture}
That etymology is not a curiosity here; it tells you the register.

\textbf{La administración} is what a notice calls itself when it wants to sound like an institution rather than a person. A letter from it will not say \emph{le pedimos} — it will say \emph{se requiere}, and it will ask for \textbf{documentación} rather than papers.

So when you see the word on a form, expect the impersonal grammar that comes with it. The vocabulary and the sentence shape travel together, and recognising one tells you to expect the other.
\end{culture}

\subsection*{Your turn}
Say these aloud:

\begin{itemize}
\raggedright
  \item \emph{la administración}
  \item she works in an office, in the public sector
  \item what \emph{minister} originally meant, and what \emph{magister} meant
\end{itemize}

\begin{tcolorbox}[breakable,colback=teal!4,colframe=teal!35!black,title={Before you move on}]
The administration? (\textbf{La administración}.) What does \emph{minus} mean? (\textbf{Less}.) And the word built on \emph{more}? (\emph{\textbf{Magister}}, behind \emph{maestro}.)
\end{tcolorbox}
\section[el envío]{el envío — put on the road}

\label{lesson:ES-C442-envio}



\begin{quote}
Say \textbf{enviar} and \textbf{la carta}. Say \textbf{la administración}. You have the verb for sending. Name the thing that is on its way.
\end{quote}

\subsection*{You'll want to know: el envío}
\textbf{el envío} — the shipment, the dispatch, the consignment.

\emph{Gastos de envío} — shipping costs, the line on every order page. \emph{Su envío está en camino} — your parcel is on its way. \emph{El envío se ha retrasado} uses the \textbf{retraso} family you already have.

The accent sits on the \textbf{í} and breaks it away from the \textbf{o}: \emph{en-VÍ-o}, three syllables.

\begin{cousinweb}
\textbf{enviar} is Latin \textbf{inviāre} — \textbf{in-} onto, plus \textbf{via}, a road.

To send something was \textbf{to put it on the road}, and the road is still sitting inside the word. \emph{El envío} is what has been put on it.

\noindent
\begin{tabularx}{\linewidth}{@{}>{\raggedright\arraybackslash}X>{\raggedright\arraybackslash}X@{}}
\toprule
\textbf{word} & \textbf{the road in it} \\
\midrule
enviar & put onto the way \\
el envío & the thing on the way \\
\bottomrule
\end{tabularx}

English has \emph{via} unchanged as a preposition, and the same \emph{via} is inside \textbf{obvious} — \emph{ob viam}, lying in the way, so plain you trip over it.
\end{cousinweb}

\begin{culture}
\textbf{Envío} is the second word in two chapters built by putting a plain ending on a verb you already own: \emph{prestar} gave \textbf{préstamo}, \emph{enviar} gives \textbf{envío}.

That is a pattern worth using deliberately. When you know a Spanish verb and need the noun for its result, the noun very often exists and is short — and you can usually recognise it even when you could not have produced it.

The administración will use the noun where a friend would use the verb. A friend says \emph{te lo envío}; the tracking page says \emph{su envío}. Same event, and the choice of word tells you who is writing.
\end{culture}

\subsection*{Your turn}
Say these aloud:

\begin{itemize}
\raggedright
  \item \emph{el envío}
  \item the shipping costs are not included
  \item what the word \emph{via} is doing inside \emph{enviar}
\end{itemize}

\begin{tcolorbox}[breakable,colback=teal!4,colframe=teal!35!black,title={Before you move on}]
The shipment? (\textbf{El envío}.) Which Latin word is hidden in it? (\emph{\textbf{Via}}, a road.) How many syllables? (\textbf{Three — en-VÍ-o}.)
\end{tcolorbox}
\section[concertar]{concertar — agreeing, by way of a fight}

\label{lesson:ES-C442-concertar}



\begin{quote}
Say \textbf{la cita} and \textbf{la hora}. Say \textbf{el envío}. You know the appointment and you know the time. Name the act of fixing them together.
\end{quote}

\subsection*{You'll want to know: concertar}
\textbf{concertar} — to arrange, to set up by agreement.

\emph{Concertar una cita} — to make an appointment. It is the verb a clinic, a bank or a council uses on its website, where a friend would say \emph{quedar}.

It breaks its stem under stress, like \emph{cerrar}: \textbf{concierto}, \textbf{conciertas}, \textbf{concierta}, but \emph{concertamos}.

\begin{cousinweb}
Say the \emph{yo} form aloud: \textbf{concierto}. You have met that word already, and it meant a performance.

\textbf{el concierto} and \textbf{yo concierto} are the same word. Both are Latin \textbf{concertāre}, and here is the part that surprises people: \emph{concertāre} meant \textbf{to contend, to fight it out} — \emph{con-} together plus \emph{certāre}, to compete.

\noindent
\begin{tabularx}{\linewidth}{@{}>{\raggedright\arraybackslash}X>{\raggedright\arraybackslash}X@{}}
\toprule
\textbf{what it became} & \textbf{how the contest turned into it} \\
\midrule
\textbf{concertar} una cita & settle it by thrashing it out \\
\textbf{un concierto} & voices contending \emph{with} each other \\
\bottomrule
\end{tabularx}

The sense slid from fighting \emph{against} to striving \emph{together}, and both descendants kept a piece of it. A negotiation and an orchestra are the same activity seen from two moods.
\end{cousinweb}

\begin{culture}
English has the same pair and hides it every bit as well. A \textbf{concert} is the performance; \textbf{to concert} an effort, and \textbf{in concert with}, are the agreement sense, and \textbf{concerted} is what several people fought out in advance.

So \emph{concertar} is not a word to learn cold. You already had \textbf{el concierto}, and the only new thing is that the noun has a verb standing behind it.

That is worth doing on purpose whenever a new verb looks like a noun you know. Check whether they are the same word before you treat them as two.
\end{culture}

\subsection*{Your turn}
Say these aloud:

\begin{itemize}
\raggedright
  \item \emph{concertar}
  \item I want to make an appointment for Thursday at nine
  \item why \emph{concierto} has two meanings
\end{itemize}

\begin{tcolorbox}[breakable,colback=teal!4,colframe=teal!35!black,title={Before you move on}]
To arrange an appointment? (\textbf{Concertar una cita}.) What did \emph{certāre} mean? (\textbf{To compete, to fight}.) What is the \emph{yo} form? (\textbf{Concierto}.)
\end{tcolorbox}
\section[la reunión]{la reunión — made one again}

\label{lesson:ES-C442-reunion}



\begin{quote}
Say \textbf{la sala} and \textbf{celebrar}. Say \textbf{concertar}. You can fix a time and you know what a notice calls holding an event. Name the event itself.
\end{quote}

\subsection*{You'll want to know: la reunión}
\textbf{la reunión} — the meeting.

\emph{Una reunión de vecinos} is a residents' meeting, the one whose notice goes up in the entrance hall. \emph{La sala de reuniones} is the meeting room. The verb is \textbf{reunirse}: \emph{nos reunimos el jueves} — we are meeting on Thursday.

\begin{cousinweb}
\textbf{re-} again, plus \textbf{unir} from Latin \textbf{ūnus}, one, plus \textbf{-ión}.

A \emph{reunión} is \textbf{a making-one-again}. The word assumes the people were together before and have scattered, which is why \emph{reunirse} also covers a family gathering and why English \textbf{reunion} sounds warmer than \textbf{meeting}.

Spanish uses the one word for both. \emph{Una reunión de trabajo} and \emph{una reunión familiar} are the same noun, and only the adjective says which.
\end{cousinweb}

\begin{culture}
Put this next to \textbf{celebrarse}, which you met on a notice meaning \emph{is held}.

\begin{quote}
La reunión de vecinos se \textbf{celebrará} el jueves.
\end{quote}

Both halves of that sentence are about a crowd. \textbf{Celebrar} is built on \emph{celeber}, thronged; \textbf{reunión} is built on making people one again. The sentence says, twice over, that people are being brought together — and neither word mentions anybody doing the bringing.

That doubling is ordinary in official Spanish. Once you can hear it, a notice stops being a wall of formal vocabulary and becomes two or three plain ideas said in their dressed-up clothes.
\end{culture}

\subsection*{Your turn}
Say these aloud:

\begin{itemize}
\raggedright
  \item \emph{la reunión}
  \item we are meeting on Thursday in the meeting room
  \item I want to arrange a meeting
\end{itemize}

\begin{tcolorbox}[breakable,colback=teal!4,colframe=teal!35!black,title={Before you move on}]
The meeting? (\textbf{La reunión}.) Which number is inside it? (\emph{\textbf{Ūnus}}, one.) What does the \textbf{re-} add? (\textbf{Again}.)
\end{tcolorbox}
\section[la pantalla]{la pantalla — the thing you put in front of a light}

\label{lesson:ES-C442-pantalla}



\begin{quote}
Say \textbf{el ordenador} and \textbf{el móvil}. Say \textbf{la reunión}. Both devices have the same part, and a meeting held remotely happens on it.
\end{quote}

\subsection*{You'll want to know: la pantalla}
\textbf{la pantalla} — the screen.

\emph{La pantalla del ordenador}, \emph{la pantalla del móvil}. \emph{Compartir pantalla} — to share your screen. \emph{En pantalla} — on screen. It also still means a lampshade, and a cinema screen: \emph{la gran pantalla} is the big screen.

The \textbf{ll} is one sound, not two: \emph{pan-TA-ya}.

\begin{culture}
The oldest sense is the useful one. A \emph{pantalla} was \textbf{a shade put in front of a flame} so the light did not hit your eyes directly.

That is one idea, and every later use is the same idea moved:

\noindent
\begin{tabularx}{\linewidth}{@{}>{\raggedright\arraybackslash}X>{\raggedright\arraybackslash}X@{}}
\toprule
\textbf{what it shades} & \textbf{the sense} \\
\midrule
a candle & the lampshade \\
a projector's beam & the cinema screen \\
a lit device & the phone or computer screen \\
\bottomrule
\end{tabularx}

Nothing about a modern screen shades anything, but the word arrived by way of things that did. Spanish did not borrow a new word for the new object; it moved an old one, and the chain is short enough to hold in your head.
\end{culture}

\begin{cousinweb}
It does not come apart, and its history is unsettled. Dictionaries reach for a crossing of two older words and do not commit further; no derivation has carried the field.

That makes three now — \textbf{el andén}, \textbf{tirar} and \textbf{la pantalla} — and it is the same honest answer each time. You met the piece-test in the chapter on the sports centre: take a piece off and see whether the rest stands. These words fail that test at the first step, because there is no piece to take off, and no amount of staring at modern Spanish supplies one.

What you \emph{can} do with a word like this is what the block above did: follow the \textbf{meanings} instead of the parts. \emph{Pantalla} has no recoverable ancestry and a perfectly clear career, and the career is what makes it memorable.
\end{cousinweb}

\subsection*{Your turn}
Say these aloud:

\begin{itemize}
\raggedright
  \item \emph{la pantalla}
  \item the meeting is on screen, I will share my screen
  \item what a \emph{pantalla} was before it was a screen
\end{itemize}

\begin{tcolorbox}[breakable,colback=teal!4,colframe=teal!35!black,title={Before you move on}]
The screen? (\textbf{La pantalla}.) What did it mean first? (\textbf{A shade in front of a light}.) Does the word come apart? (\textbf{No, and its origin is unsettled}.)
\end{tcolorbox}
\section[(administración, envío, concertar …]{(administración, envío, concertar, reunión, pantalla) — from cold}

\label{lesson:ES-R442-repaso-gestion}



\begin{quote}
Close the lessons before this one. Say all five: \emph{officialdom}, \emph{the shipment}, \emph{to arrange an appointment}, \emph{the meeting}, \emph{the screen}.
\end{quote}

\begin{grammarlens}[title={Grammar Lens: two routes into a word, and when each one runs out}]
Four of this chapter's five words come apart cleanly.

\noindent
\begin{tabularx}{\linewidth}{@{}>{\raggedright\arraybackslash}X>{\raggedright\arraybackslash}X@{}}
\toprule
\textbf{word} & \textbf{the part that carries it} \\
\midrule
administración & \emph{minus}, less — the one who serves \\
reunión & \emph{ūnus}, one — made one again \\
\textbf{pantalla} & \textbf{nothing; follow the meanings instead} \\
\bottomrule
\end{tabularx}

For \textbf{pantalla} the parts route is closed, and a different one is open: a shade in front of a flame, then in front of a projector, then in front of a lit device. One idea moved three times.

Both routes make a word stick. What does not work is insisting on the first when only the second is available, because that is how invented etymologies get made.
\end{grammarlens}

\subsection*{Your turn}
Say these aloud:

\begin{itemize}
\raggedright
  \item I want to arrange a meeting for Thursday
  \item the shipping costs are not included
  \item \emph{concertar} and \emph{un concierto}, and why they are one word
\end{itemize}

\begin{grammarlens}[title={What you've built}]
\textbf{You can now deal with an institution rather than a person} — book the appointment, read what the office wants, track what it sent you, and do it on a screen. That is the half of A2 reading that is not about friends.

The register travels with the vocabulary. \textbf{La administración} does not write \emph{le pedimos}; it writes \emph{se requiere}. When one of these words appears, expect the impersonal grammar that comes with it.
\end{grammarlens}

\begin{tcolorbox}[breakable,colback=teal!4,colframe=teal!35!black,title={Before you move on}]
The meeting? (\textbf{La reunión}.) The screen? (\textbf{La pantalla}.) Officialdom? (\textbf{La administración}.) Which of the three has no recoverable origin? (\emph{\textbf{Pantalla}}.)
\end{tcolorbox}
\section[Practice]{(una cita con la administración) — read it, then do it yourself}

\label{lesson:ES-C442-sintesis-gestion}



\begin{quote}
Say \emph{la administración} and \emph{concertar}. The page you are about to read uses both, and it never once says \emph{you}.
\end{quote}

\subsection*{Your turn}
Read it once straight through, then answer.

\textbf{en la web de la administración}

\begin{quote}
Para \textbf{concertar} una cita, seleccione el trámite en la \textbf{pantalla} de inicio. Se requiere documentación en vigor. La cita se confirmará por correo. Si no puede acudir, se ruega anular con 24 horas de antelación.
\end{quote}

Say these aloud:

\begin{itemize}
\raggedright
  \item what you have to do first
  \item what you must bring
  \item what to do if you cannot come
\end{itemize}

\begin{culture}
Count the people in that notice. There are none.

\textbf{Se requiere}, \textbf{se confirmará}, \textbf{se ruega} — three verbs, no subject. The one exception is \emph{seleccione}, an instruction aimed at whoever is reading, and even that is the formal form that avoids saying \emph{tú}.

This is the same \textbf{se} you met on \emph{la junta se celebra} and \emph{la oficina ha sido trasladada}: it takes the person out and leaves the obligation attached to the procedure. A public-sector page is written this way from top to bottom, and once you expect it, the sentences get shorter — \emph{se requiere documentación} means no more than \emph{bring your papers}.

The tracking page for your \textbf{envío} and the notice about the \textbf{reunión} use exactly the same grammar. One register, three settings.
\end{culture}

\begin{tcolorbox}[breakable,colback=teal!4,colframe=teal!35!black,title={Before you move on}]
Now your turn, out loud and then in writing. Five sentences: say you want to make an appointment, say what the screen asked for, say what papers you took, say the meeting was moved, and say your parcel has not come.

Then say the first one again the way the website would say it, with no \emph{you} in it at all.
\end{tcolorbox}
